% i18n_qa_mixed_zh_en_paper.tex
% Chinese paper with English abstract
% Expected detect_language result: "zh" (CJK heuristic via xeCJK)
% Deliberate errors: Western punctuation in Chinese text (commas, periods)
\documentclass[12pt,a4paper]{article}
\usepackage{xeCJK}
\setCJKmainfont{SimSun}
\setCJKsansfont{SimHei}
\setCJKmonofont{FangSong}
\usepackage{amsmath,amssymb,amsthm}
\usepackage{graphicx}
\usepackage{booktabs}
\usepackage{hyperref}
\usepackage{algorithm}
\usepackage{algpseudocode}
\usepackage{tikz}
\usepackage{pgfplots}
\pgfplotsset{compat=1.18}
\usepackage{enumitem}
\usepackage{float}

\newtheorem{theorem}{定理}[section]
\newtheorem{lemma}[theorem]{引理}
\newtheorem{proposition}[theorem]{命题}
\newtheorem{corollary}[theorem]{推论}
\theoremstyle{definition}
\newtheorem{definition}{定义}[section]
\newtheorem{remark}{注}[section]

\title{基于注意力机制的多模态特征融合方法研究\\
  {\large 及其在医学图像分析中的应用}}
\author{王建华\textsuperscript{1}\quad 李明辉\textsuperscript{2}\quad
  张晓燕\textsuperscript{1}\\[6pt]
  \textsuperscript{1}清华大学计算机科学与技术系\\
  \textsuperscript{2}北京大学智能科学系}
\date{2026年3月}

\begin{document}
\maketitle

\begin{abstract}
\noindent\textbf{Abstract.}
This paper proposes a novel multi-modal feature fusion framework based on
cross-attention mechanisms for medical image analysis. Our approach, termed
Cross-Modal Attention Fusion Network (CMAFN), effectively integrates
information from CT scans, MRI images, and clinical text reports to improve
diagnostic accuracy for pulmonary diseases. Unlike existing late-fusion
strategies that simply concatenate features from different modalities,
CMAFN learns modality-specific attention weights that dynamically adjust
the contribution of each input source based on the diagnostic context.
We evaluate our method on two large-scale datasets: the ChestX-ray14
benchmark containing 112,120 frontal-view chest X-rays and a proprietary
multi-modal dataset from Beijing Union Medical College Hospital comprising
15,432 patient records with paired CT, MRI, and radiology reports. Our
experimental results demonstrate that CMAFN achieves state-of-the-art
performance with an AUC of 0.943 on ChestX-ray14 and 0.961 on the
multi-modal dataset, outperforming single-modality baselines by 4.2\%
and 6.8\% respectively. Ablation studies confirm the effectiveness of
the cross-attention mechanism and the importance of incorporating textual
clinical information alongside imaging data.

\medskip
\noindent\textbf{Keywords:} multi-modal learning, attention mechanism,
medical image analysis, feature fusion, deep learning.
\end{abstract}


\section{引言}

% ERROR: Western comma used instead of Chinese comma (，)
医学图像分析是计算机视觉领域的一个重要研究方向,近年来深度学习技术的
快速发展为该领域带来了革命性的变化。传统的医学诊断主要依赖于放射科医生
对单一模态影像的视觉判读,这种方式不仅耗时耗力,而且容易受到医生主观
经验的影响。

% ERROR: Western period used instead of Chinese period (。)
随着医疗信息化的推进,患者的诊疗数据日益丰富,包括CT扫描、MRI成像、
X光片、超声图像以及临床文本报告等多种模态的信息.如何有效地融合这些
多模态信息以提高诊断准确率,是当前医学人工智能研究的核心问题之一。

在本文中，我们提出了一种基于交叉注意力机制的多模态特征融合网络
（Cross-Modal Attention Fusion Network，简称CMAFN），该方法能够
自适应地学习不同模态之间的关联关系，并根据诊断任务动态调整各模态的
贡献权重。

本文的主要贡献包括以下几个方面：
\begin{enumerate}[label=(\arabic*)]
  \item 提出了一种新颖的交叉注意力融合机制，能够捕捉不同模态特征之间
    的细粒度交互关系；
  \item 设计了一个端到端的多模态学习框架，支持影像和文本两类输入的
    联合训练；
  \item 在两个大规模数据集上验证了方法的有效性，取得了领先的诊断性能；
  \item 通过消融实验和可视化分析，深入探讨了各组件对最终性能的贡献。
\end{enumerate}


\section{相关工作}

\subsection{医学图像分析}

% ERROR: Western comma in Chinese context
深度卷积神经网络在医学图像分析中的应用始于2012年Krizhevsky等人
\cite{krizhevsky2012imagenet}提出的AlexNet,此后ResNet
\cite{he2016deep}、DenseNet \cite{huang2017densely}等经典网络结构被
广泛应用于各类医学影像任务。在胸部X光片分析方面,Wang等人
\cite{wang2017chestx}构建了大规模的ChestX-ray14数据集并提出了基线模型。

Rajpurkar等人\cite{rajpurkar2017chexnet}提出的CheXNet模型在14种胸部
疾病的检测任务上达到了放射科医生水平的性能。然而，这些方法大多仅利用
单一模态的影像信息，忽略了临床文本报告中包含的丰富语义信息。

\subsection{多模态学习}

多模态学习旨在利用来自不同数据源的互补信息来提升模型性能。根据融合
策略的不同，现有方法可分为三类：

\begin{definition}[多模态融合策略]
设 $\{M_1, M_2, \ldots, M_K\}$ 为 $K$ 种不同模态的特征表示，
多模态融合函数 $\mathcal{F}$ 定义为：
\begin{equation}\label{eq:fusion-def}
  z = \mathcal{F}(M_1, M_2, \ldots, M_K; \theta_f),
\end{equation}
其中 $z$ 为融合后的表示，$\theta_f$ 为融合参数。
\end{definition}

% ERROR: Western period in Chinese sentence
早期融合（Early Fusion）在输入层直接拼接各模态的原始特征.
晚期融合（Late Fusion）分别对各模态进行独立编码后在决策层进行融合.
中间融合（Intermediate Fusion）在网络的中间层进行特征交互。

\subsection{注意力机制}

注意力机制最初由Bahdanau等人\cite{bahdanau2015neural}在机器翻译任务中
提出，后来Vaswani等人\cite{vaswani2017attention}提出的Transformer架构
将自注意力机制推广到了更广泛的应用场景。在多模态学习中，交叉注意力
（Cross-Attention）机制被用于建模不同模态之间的对应关系：
\begin{equation}\label{eq:cross-attn}
  \text{CrossAttn}(Q_i, K_j, V_j) =
    \text{softmax}\!\left(\frac{Q_i K_j^\top}{\sqrt{d_k}}\right) V_j,
\end{equation}
其中 $Q_i$ 来自模态 $i$，$K_j$ 和 $V_j$ 来自模态 $j$，
$d_k$ 为键向量的维度。


\section{方法}

\subsection{整体框架}

CMAFN由三个主要模块组成：模态编码器（Modality Encoder）、交叉注意力
融合模块（Cross-Attention Fusion Module）和分类头（Classification
Head）。设输入包含影像 $I \in \mathbb{R}^{H \times W \times C}$ 和
文本报告 $T = (t_1, t_2, \ldots, t_L)$，其中 $L$ 为文本长度。

\begin{definition}[模态编码器]
影像编码器 $E_v$ 和文本编码器 $E_t$ 分别将输入映射到统一的特征空间：
\begin{align}
  F_v &= E_v(I; \theta_v) \in \mathbb{R}^{N_v \times d}, \label{eq:vis-enc} \\
  F_t &= E_t(T; \theta_t) \in \mathbb{R}^{N_t \times d}, \label{eq:txt-enc}
\end{align}
其中 $N_v$ 为影像特征的空间位置数，$N_t$ 为文本标记数，
$d$ 为特征维度。
\end{definition}

\subsection{交叉注意力融合模块}

融合模块的核心是双向交叉注意力层，其数学形式如下：
\begin{align}
  \hat{F}_v &= F_v + \text{CrossAttn}(F_v W_Q^v, F_t W_K^t, F_t W_V^t),
    \label{eq:v2t-attn} \\
  \hat{F}_t &= F_t + \text{CrossAttn}(F_t W_Q^t, F_v W_K^v, F_v W_V^v),
    \label{eq:t2v-attn}
\end{align}
其中 $W_Q^v, W_K^t, W_V^t, W_Q^t, W_K^v, W_V^v \in \mathbb{R}^{d \times d_k}$
为可学习的投影矩阵。

\begin{lemma}[融合表示的信息保持性]
设 $H(X)$ 表示随机变量 $X$ 的信息熵。在交叉注意力融合下：
\begin{equation}
  H(\hat{F}_v, \hat{F}_t) \geq \max\{H(F_v), H(F_t)\},
\end{equation}
即融合后的联合表示至少保留了各单模态表示的信息量。
\end{lemma}

\subsection{多尺度特征提取}

为了捕捉不同尺度的病变特征，我们在影像编码器中引入了特征金字塔
网络（FPN）：
\begin{equation}\label{eq:fpn}
  P_l = \text{Conv}_{1\times 1}(C_l) + \text{Upsample}(P_{l+1}),
    \quad l = 3, 4, 5,
\end{equation}
其中 $C_l$ 为ResNet第 $l$ 层的特征图，$P_l$ 为对应的金字塔层特征。

\subsection{损失函数}

% ERROR: Western comma in Chinese text
总损失函数由分类损失和对齐损失两部分组成,定义如下：
\begin{equation}\label{eq:total-loss}
  \mathcal{L} = \mathcal{L}_{\text{cls}} + \alpha \mathcal{L}_{\text{align}},
\end{equation}
其中分类损失采用加权二元交叉熵：
\begin{equation}\label{eq:cls-loss}
  \mathcal{L}_{\text{cls}} = -\sum_{c=1}^{C} \left[
    w_c y_c \log(\hat{y}_c) + (1-y_c) \log(1-\hat{y}_c)
  \right],
\end{equation}
对齐损失鼓励语义相似的跨模态特征在嵌入空间中靠近：
\begin{equation}\label{eq:align-loss}
  \mathcal{L}_{\text{align}} = \sum_{i,j} \max\left(0,
    \|f_v^{(i)} - f_t^{(j)}\|_2 - m + s_{ij}\right),
\end{equation}
其中 $m$ 为间隔参数，$s_{ij}$ 表示影像区域 $i$ 和文本片段 $j$
之间的语义相似度。


\section{实验}

\subsection{数据集}

我们在两个数据集上进行实验验证：

\begin{table}[H]
\centering
\caption{实验数据集统计信息}\label{tab:datasets}
\begin{tabular}{@{}lrrrr@{}}
\toprule
数据集 & 样本数 & 疾病类别 & 模态 & 来源 \\
\midrule
ChestX-ray14 & 112,120 & 14 & X光 & NIH \\
BUMCH-MM & 15,432 & 8 & CT+MRI+文本 & 北京协和医院 \\
\bottomrule
\end{tabular}
\end{table}

\subsection{实现细节}

影像编码器采用在ImageNet上预训练的ResNet-50 \cite{he2016deep}作为
骨干网络。文本编码器采用预训练的Chinese-BERT-wwm
\cite{cui2020revisiting}。交叉注意力模块使用8头注意力，特征维度
$d = 512$，键/值维度 $d_k = 64$。

\begin{table}[H]
\centering
\caption{训练超参数设置}\label{tab:hyperparams}
\begin{tabular}{@{}lr@{}}
\toprule
超参数 & 值 \\
\midrule
学习率 & $1 \times 10^{-4}$ \\
批量大小 & 32 \\
优化器 & AdamW \\
权重衰减 & $1 \times 10^{-5}$ \\
训练轮数 & 100 \\
对齐损失权重 $\alpha$ & 0.1 \\
间隔参数 $m$ & 0.5 \\
\bottomrule
\end{tabular}
\end{table}

\subsection{主要结果}

\begin{table}[H]
\centering
\caption{ChestX-ray14数据集上各方法的AUC比较}\label{tab:main-results}
\begin{tabular}{@{}lcccc@{}}
\toprule
方法 & 肺不张 & 心脏肥大 & 渗出 & 平均AUC \\
\midrule
ResNet-50 \cite{he2016deep}  & 0.816 & 0.904 & 0.884 & 0.901 \\
DenseNet-121 \cite{huang2017densely} & 0.826 & 0.912 & 0.891 & 0.910 \\
CheXNet \cite{rajpurkar2017chexnet} & 0.832 & 0.918 & 0.897 & 0.917 \\
Late Fusion & 0.845 & 0.925 & 0.905 & 0.928 \\
Bilinear Pooling & 0.851 & 0.929 & 0.911 & 0.933 \\
\textbf{CMAFN (ours)} & \textbf{0.872} & \textbf{0.947} & \textbf{0.932} & \textbf{0.943} \\
\bottomrule
\end{tabular}
\end{table}

\subsection{消融实验}

% ERROR: Western period in Chinese sentence
为了验证各模块的有效性,我们进行了详细的消融实验.具体结果如表
\ref{tab:ablation}所示。

\begin{table}[H]
\centering
\caption{消融实验结果（BUMCH-MM数据集）}\label{tab:ablation}
\begin{tabular}{@{}lccc@{}}
\toprule
配置 & AUC & Precision & Recall \\
\midrule
仅影像 & 0.894 & 0.856 & 0.831 \\
仅文本 & 0.867 & 0.823 & 0.798 \\
晚期融合 & 0.931 & 0.892 & 0.876 \\
无交叉注意力 & 0.938 & 0.901 & 0.883 \\
无FPN & 0.947 & 0.912 & 0.895 \\
无对齐损失 & 0.953 & 0.918 & 0.902 \\
\textbf{CMAFN完整模型} & \textbf{0.961} & \textbf{0.928} & \textbf{0.915} \\
\bottomrule
\end{tabular}
\end{table}

\subsection{可视化分析}

图\ref{fig:attention-vis}展示了交叉注意力权重的可视化结果。可以观察到，
当文本报告中提到特定病变区域时，模型能够准确地将注意力集中到对应的
影像区域。

\begin{figure}[H]
\centering
\begin{tikzpicture}
\draw[fill=gray!20] (0,0) rectangle (4,4);
\node at (2,2) {CT影像};
\draw[->,red,ultra thick] (4.2,3) -- (5.8,3.5);
\draw[->,red,ultra thick] (4.2,1.5) -- (5.8,1);
\draw[fill=blue!10] (6,0) rectangle (10,4);
\node[text width=3.5cm,align=center] at (8,3.5) {右上肺叶};
\node[text width=3.5cm,align=center] at (8,2.5) {磨玻璃影};
\node[text width=3.5cm,align=center] at (8,1.5) {结节状病变};
\node[text width=3.5cm,align=center] at (8,0.5) {考虑肺癌};
\draw[red,ultra thick,dashed] (1,2.5) circle (0.8);
\node[red,font=\small] at (1,1.2) {高注意力区域};
\end{tikzpicture}
\caption{交叉注意力可视化示例：影像区域与文本描述的对应关系}\label{fig:attention-vis}
\end{figure}


\section{讨论}

% ERROR: Western comma mixed with Chinese text
实验结果表明,CMAFN在两个数据集上均取得了优异的性能。与单模态方法
相比,多模态融合带来了显著的性能提升,这证实了不同模态之间存在互补
信息。

\begin{remark}
值得注意的是，交叉注意力机制不仅提高了分类性能，还提供了良好的
可解释性。通过可视化注意力权重，临床医生可以直观地理解模型的
诊断依据，这对于医学人工智能系统的临床应用至关重要。
\end{remark}

\subsection{局限性}

本研究存在以下局限性：首先，BUMCH-MM数据集来自单一医疗机构，模型在
跨机构泛化方面的表现有待进一步验证。其次，当前方法假设所有模态数据
都是完整的，而实际临床场景中经常存在模态缺失的情况。最后，模型的
计算复杂度与模态数量呈平方关系，可扩展性需要进一步优化。


\section{结论}

% ERROR: Western period ending Chinese sentence
本文提出了CMAFN多模态融合框架,有效地结合了医学影像和临床文本信息.
实验结果验证了该方法的有效性和优越性。未来工作将围绕模态缺失鲁棒性、
跨机构泛化以及三维影像处理等方向展开深入研究。


\begin{thebibliography}{99}

\bibitem{krizhevsky2012imagenet}
A.~Krizhevsky, I.~Sutskever, and G.~E. Hinton.
\newblock {ImageNet} classification with deep convolutional neural networks.
\newblock In \textit{Advances in Neural Information Processing Systems
  (NeurIPS)}, pages 1097--1105, 2012.

\bibitem{he2016deep}
K.~He, X.~Zhang, S.~Ren, and J.~Sun.
\newblock Deep residual learning for image recognition.
\newblock In \textit{IEEE Conference on Computer Vision and Pattern Recognition
  (CVPR)}, pages 770--778, 2016.

\bibitem{huang2017densely}
G.~Huang, Z.~Liu, L.~van~der~Maaten, and K.~Q. Weinberger.
\newblock Densely connected convolutional networks.
\newblock In \textit{IEEE Conference on Computer Vision and Pattern Recognition
  (CVPR)}, pages 4700--4708, 2017.

\bibitem{wang2017chestx}
X.~Wang, Y.~Peng, L.~Lu, Z.~Lu, M.~Bagheri, and R.~M. Summers.
\newblock {ChestX-ray8}: Hospital-scale chest X-ray database and benchmarks.
\newblock In \textit{IEEE Conference on Computer Vision and Pattern Recognition
  (CVPR)}, pages 2097--2106, 2017.

\bibitem{rajpurkar2017chexnet}
P.~Rajpurkar, J.~Irvin, K.~Zhu, et~al.
\newblock {CheXNet}: Radiologist-level pneumonia detection on chest X-rays with
  deep learning.
\newblock \textit{arXiv preprint arXiv:1711.05225}, 2017.

\bibitem{vaswani2017attention}
A.~Vaswani, N.~Shazeer, N.~Parmar, et~al.
\newblock Attention is all you need.
\newblock In \textit{Advances in Neural Information Processing Systems
  (NeurIPS)}, pages 5998--6008, 2017.

\bibitem{bahdanau2015neural}
D.~Bahdanau, K.~Cho, and Y.~Bengio.
\newblock Neural machine translation by jointly learning to align and translate.
\newblock In \textit{International Conference on Learning Representations
  (ICLR)}, 2015.

\bibitem{cui2020revisiting}
Y.~Cui, W.~Che, T.~Liu, B.~Qin, and Z.~Yang.
\newblock Revisiting pre-trained models for {Chinese} natural language
  processing.
\newblock In \textit{Findings of EMNLP}, pages 657--668, 2020.

\bibitem{lu2019vilbert}
J.~Lu, D.~Batra, D.~Parikh, and S.~Lee.
\newblock {ViLBERT}: Pretraining task-agnostic visiolinguistic representations
  for vision-and-language tasks.
\newblock In \textit{Advances in Neural Information Processing Systems
  (NeurIPS)}, pages 13--23, 2019.

\bibitem{chen2020uniter}
Y.~Chen, L.~Li, L.~Yu, et~al.
\newblock {UNITER}: Universal image-text representation learning.
\newblock In \textit{European Conference on Computer Vision (ECCV)},
  pages 104--120, 2020.

\end{thebibliography}

\end{document}
